% from ega2-1-fr.tex
\oldpage[II]{38}
section
$(z|D_+(f))(\alpha_d(f^k)|D_+(f))$ de $\mathcal{F}((n+k)d)$, soit nulle,
alors, en raison de la remarque précédente, on a aussi $z|D_+(f)=0$. Cela
montre que l'on définit un $S_{(f)}$-homomorphisme
$\beta_f:M_{(f)}\to\Gamma(D_+(f),\mathcal{F})$ en faisant correspondre à
l'élément $z/f^n$ la section
$(z|D_+(f))(\alpha_d(f^n)|D_+(f))^{-1}$ de $\mathcal{F}$ au-dessus de
$D_+(f)$. On vérifie en outre aussitôt que le diagramme
\[
\label{II.2.6.4.1-fr}
  \xymatrix{
    M_{(f)}\ar[r]^{\beta_f}\ar[d]
    & \Gamma(D_+(f),\mathcal{F})\ar[d]\\
    M_{(fg)}\ar[r]_{\beta_{fg}}
    & \Gamma(D_+(fg),\mathcal{F})
  }
\tag{2.6.4.1}
\]
est commutatif pour $g\in S_e$ ($e>0$). Si on se rappelle que $M_{(f)}$
s'identifie canoniquement à $\Gamma(D_+(f),\widetilde{M})$ et que les
$D_+(f)$ forment une base de la topologie de $X$
(\hyperref[II.2.3.4-fr]{2.3.4}), on voit que les $\beta_f$ proviennent d'un
unique homomorphisme canonique de $\mathcal{O}_X$-Modules
\[
\label{II.2.6.4.2-fr}
  \beta:\widetilde{\Gamma_*(\mathcal{F})}
  \longrightarrow\mathcal{F}
\tag{2.6.4.2}
\]
(aussi noté $\beta_{\mathcal{F}}$) qui est évidemment fonctoriel.
\end{env}

\begin{proposition}[2.6.5]
\label{II.2.6.5-fr}
Soient $M$ un $S$-module gradué, $\mathcal{F}$ un $\mathcal{O}_X$-Module~;
alors les homomorphismes composés
\[
\label{II.2.6.5.1-fr}
  \widetilde{M}\xrightarrow{\widetilde{\alpha}}
  \widetilde{\Gamma_*(\widetilde{M})}
  \xrightarrow{\beta}\widetilde{M}
\tag{2.6.5.1}
\]
\[
\label{II.2.6.5.2-fr}
  \Gamma_*(\mathcal{F})\xrightarrow{\alpha}
  \Gamma_*(\widetilde{\Gamma_*(\mathcal{F})})
  \xrightarrow{\Gamma_*(\beta)}\Gamma_*(\mathcal{F})
\tag{2.6.5.2}
\]
sont les isomorphismes identiques.
\end{proposition}

\begin{proof}
La vérification de (\hyperref[II.2.6.5.1-fr]{2.6.5.1}) est locale~: dans un
ouvert $D_+(f)$, elle résulte aussitôt des définitions et de ce que $\beta$,
appliqué à des faisceaux quasi-cohérents, est déterminé par son action sur les
sections au-dessus de $D_+(f)$ (\hyperref[I.1.3.8-fr]{I, 1.3.8}). La
vérification de (\hyperref[II.2.6.5.2-fr]{2.6.5.2}) se fait pour chaque degré
séparément~: si on pose $M=\Gamma_*(\mathcal{F})$, on a
$M_n=\Gamma(X,\mathcal{F}(n))$ et
$(\Gamma_*(\widetilde{M}))_n=\Gamma(X,\widetilde{M}(n))
=\Gamma(X,\widetilde{M(n)})$. Or, si $f\in S_1$ et $z\in M_n$,
$\alpha_n^f(z)$ est l'élément $z/1$ de $(M(n))_{(f)}$, égal à
$(f/1)^n(z/f^n)$~; il lui correspond par $\beta_f$ la section
\[
  ((\alpha_1(f))^n|D_+(f))
  ((z|D_+(f))((\alpha_1(f))^n|D_+(f))^{-1})
\]
au-dessus de $D_+(f)$, c'est-à-dire la restriction de $z$ à $D_+(f)$, ce qui
vérifie (\hyperref[II.2.6.5.2-fr]{2.6.5.2}).
\end{proof}

\subsection{Conditions de finitude.}
\label{subsection:II.2.7-fr}

\begin{proposition}[2.7.1]
\label{II.2.7.1-fr}
\begin{enumerate}
  \item[\rm(i)] Si $S$ est un anneau gradué noethérien,
    $X=\operatorname{Proj}(S)$ est un schéma noethérien.
  \item[\rm(ii)] Si $S$ est une $A$-algèbre graduée de type fini,
    $X=\operatorname{Proj}(S)$ est un schéma de type fini au-dessus de
    $Y=\operatorname{Spec}(A)$.
\end{enumerate}
\end{proposition}

